\definicion{Modelo cliente-servidor}{Arquitectura web donde el cliente (navegador) solicita recursos y el servidor los aloja y responde, complementando lenguajes de cliente (HTML, CSS, JS) y de servidor.}
\subtituloA{Componentes del navegador}

\lista{
  \textbf{Componentes:} Interfaz, buscador, renderizado, comunicaciones, intérprete JS, parser XML, visualización y persistencia.
}
\nota{Para el examen}{
Recordar que JavaScript es la capa de comportamiento de la web y que el sitio debe funcionar sin la intervención directa del servidor para la presentación inicial.
}
